% Experiment 19: File Allocation Strategy Simulator

\section{Experiment 19: File Allocation Strategy Simulator}

\subsection{Problem Statement}
The objective is to implement \texttt{filealloc} to simulate disk block allocation.

The tool must:
\begin{itemize}
    \item Read total blocks, free block list, and file requirements from stdin.
    \item Simulate \textbf{Contiguous} (first fit run), \textbf{Linked} (chain), and \textbf{Indexed} (index + data) allocation.
    \item Output the block mapping for each file or \verb|FAIL|.
    \item Reset the disk state between algorithms.
\end{itemize}

\subsection{Source Code}
\begin{lstlisting}[language=C, caption={filealloc.c}]
// ########## File Allocation Strategy Simulation ##########
// filealloc

#include <stdio.h>
#include <stdlib.h>
#include <string.h>
#include <stdbool.h>

#define E_INPUT "E_INPUT"
#define E_RANGE "E_RANGE"
#define E_DUPBLOCK "E_DUPBLOCK"

void print_error(const char *code, const char *msg) {
    fprintf(stderr, "ERROR: %s: %s\n", code, msg);
    exit(1);
}

// Comparison for sorting free blocks
int cmp_int(const void *a, const void *b) {
    return (*(int *)a - *(int *)b);
}

typedef struct {
    char name[32];
    int size;
} FileReq;

void reset_disk(bool *disk, int N, int *free_blocks, int F) {
    for (int i = 0; i < N; i++) disk[i] = true; 
    for (int i = 0; i < F; i++) {
        if (free_blocks[i] >= 0 && free_blocks[i] < N) {
            disk[free_blocks[i]] = false; 
        }
    }
}

void solve_contiguous(int N, bool *disk, FileReq *files, int M) {
    printf("ALG CONTIGUOUS\n");
    for (int i = 0; i < M; i++) {
        int best_start = -1;
        // Find first run of 'size' free blocks
        for (int start = 0; start <= N - files[i].size; start++) {
            bool fits = true;
            for (int k = 0; k < files[i].size; k++) {
                if (disk[start + k]) {
                    fits = false;
                    break;
                }
            }
            if (fits) {
                best_start = start;
                break; 
            }
        }

        if (best_start != -1) {
            printf("FILE %s -> START %d LEN %d\n", files[i].name, best_start, files[i].size);
            for (int k = 0; k < files[i].size; k++) disk[best_start + k] = true;
        } else {
            printf("FILE %s -> FAIL\n", files[i].name);
        }
    }
}

void solve_linked(int N, bool *disk, FileReq *files, int M) {
    printf("ALG LINKED\n");
    for (int i = 0; i < M; i++) {
        int needed = files[i].size;
        int allocated[needed];
        int count = 0;

        // Greedy allocation of smallest available blocks
        for (int b = 0; b < N && count < needed; b++) {
            if (!disk[b]) {
                allocated[count++] = b;
            }
        }

        if (count == needed) {
            printf("FILE %s -> CHAIN", files[i].name);
            for (int k = 0; k < count; k++) {
                printf("%s%d", (k == 0 ? " " : "->"), allocated[k]);
                disk[allocated[k]] = true;
            }
            printf("\n");
        } else {
            printf("FILE %s -> FAIL\n", files[i].name);
        }
    }
}

void solve_indexed(int N, bool *disk, FileReq *files, int M) {
    printf("ALG INDEXED\n");
    for (int i = 0; i < M; i++) {
        int needed = files[i].size + 1; 
        int allocated[needed];
        int count = 0;

        for (int b = 0; b < N && count < needed; b++) {
            if (!disk[b]) {
                allocated[count++] = b;
            }
        }

        if (count == needed) {
        // Index block is the first one allocated (smallest ID)
            int index_block = allocated[0];
            printf("FILE %s -> INDEX %d DATA", files[i].name, index_block);
            for (int k = 1; k < count; k++) {
                printf("%s%d", (k == 1 ? " " : ","), allocated[k]);
            }
            printf("\n");
        // Mark occupied
            for (int k = 0; k < count; k++) disk[allocated[k]] = true;
        } else {
            printf("FILE %s -> FAIL\n", files[i].name);
        }
    }
}

int main() {
    int N, F, M;
    
    // Read Total Blocks
    if (scanf("%d", &N) != 1) return 0;
    if (N < 1 || N > 10000) print_error(E_RANGE, "N 1..10000");

    // Read Free Blocks
    if (scanf("%d", &F) != 1) print_error(E_INPUT, "missing F");
    int *free_blocks = malloc(sizeof(int) * F);
    for (int i = 0; i < F; i++) {
        if (scanf("%d", &free_blocks[i]) != 1) print_error(E_INPUT, "bad free block");
        if (free_blocks[i] < 0 || free_blocks[i] >= N) print_error(E_RANGE, "block out of range");
    }
    
    // Sort free blocks for deterministic allocation (smallest ID first)
    qsort(free_blocks, F, sizeof(int), cmp_int);

    // Check duplicates
    for (int i = 0; i < F - 1; i++) {
        if (free_blocks[i] == free_blocks[i+1]) print_error(E_DUPBLOCK, "free block list must contain unique IDs");
    }

    // Read Files
    if (scanf("%d", &M) != 1) print_error(E_INPUT, "missing M");
    FileReq *files = malloc(sizeof(FileReq) * M);
    for (int i = 0; i < M; i++) {
        if (scanf("%s %d", files[i].name, &files[i].size) != 2) print_error(E_INPUT, "bad file spec");
        if (files[i].size < 1) print_error(E_INPUT, "bad file size");
    }

    // Simulation State
    bool *disk = malloc(sizeof(bool) * N);

    reset_disk(disk, N, free_blocks, F);
    solve_contiguous(N, disk, files, M);

    reset_disk(disk, N, free_blocks, F);
    solve_linked(N, disk, files, M);

    reset_disk(disk, N, free_blocks, F);
    solve_indexed(N, disk, files, M);

    free(free_blocks);
    free(files);
    free(disk);
    return 0;
}
\end{lstlisting}

\subsection{Sample Input and Output}

\textbf{Case 1:}
\begin{verbatim}
$ printf "10\n6\n0 1 2 5 6 9\n2\nA 2\nB 3" | ./filealloc
ALG CONTIGUOUS
FILE A -> START 0 LEN 2
FILE B -> FAIL
ALG LINKED
FILE A -> CHAIN 0->1
FILE B -> CHAIN 2->5->6
ALG INDEXED
FILE A -> INDEX 0 DATA 1,2
FILE B -> FAIL
\end{verbatim}